\documentclass[12pt]{article}

%-------------------------------------------------
% PAGE & LANGUAGE
% -------------------------------------------------

\usepackage[a4paper,margin=2cm]{geometry}

% -------------------------------------------------
% MATHEMATICS
% -------------------------------------------------

\usepackage{amsmath}
\usepackage{amssymb}
\usepackage{amsthm}
\usepackage{mathtools}
\usepackage{yhmath}

% Physics notation
\usepackage{physics}

% -------------------------------------------------
% GRAPHICS & PLOTS
% -------------------------------------------------

\usepackage{graphicx}
\usepackage{tikz}
\usepackage{pgfplots}

\pgfplotsset{compat=1.18}
\graphicspath{{./images/}}

% -------------------------------------------------
% COLORS & BOXES
% -------------------------------------------------

\usepackage[dvipsnames,svgnames,x11names]{xcolor}
\usepackage{tcolorbox}
\tcbuselibrary{skins,breakable}
\newtcolorbox{demobox}[1][]{
    enhanced,
    breakable,
    colback=black!2,
    colframe=black!70,
    boxrule=0pt,
    borderline west={2pt}{0pt}{black!70},
    sharp corners,
    left=10pt,
    right=7pt,
    top=6pt,
    bottom=6pt,
    fonttitle=\bfseries,
    title={Demostració},
    #1
}

\newtcolorbox{examplebox}[1]{
    enhanced,
    breakable,
    colback=black!1,
    colframe=black!15,
    boxrule=0.5pt,
    arc=2mm,
    borderline west={2pt}{0pt}{black!55},
    left=12pt,
    right=12pt,
    top=10pt,
    bottom=10pt,
    before skip=12pt,
    after skip=12pt,
    title={\textbf{Exemple #1}},
    fonttitle=\normalsize,
    coltitle=black,
    attach boxed title to top left={
        xshift=8pt,
        yshift=-\tcboxedtitleheight/2
    },
    boxed title style={
        colback=white,
        colframe=white,
        boxrule=0pt,
        left=3pt,
        right=3pt,
        top=1pt,
        bottom=1pt
    },
    before upper={\setlength{\parskip}{7pt}}
}

% -------------------------------------------------
% DOCUMENT FORMATTING
% -------------------------------------------------

\usepackage{enumitem}
\usepackage[expansion=false]{microtype}
\usepackage{booktabs}
\usepackage{float}
\usepackage{ragged2e}

% -------------------------------------------------
% CHEMISTRY
% -------------------------------------------------

\usepackage[version=4]{mhchem}

% -------------------------------------------------
% LINKS
% --------------------------------------------------

\usepackage[colorlinks=true,linkcolor=Black,urlcolor=Black,citecolor=Black]{hyperref}
\usepackage{tocloft}
\usepackage{cite}

% -------------------------------------------------
% CUSTOM COMMANDS
% -------------------------------------------------

\newcommand{\R}{\mathbb{R}}
\newcommand{\N}{\mathbb{N}}
\newcommand{\Z}{\mathbb{Z}}

% -------------------------------------------------
% DOCUMENT STYLE
% -------------------------------------------------

\linespread{1.25}

\setlength{\parskip}{0.8em}
\setlength{\parindent}{0pt}

\setlength{\heavyrulewidth}{0.4pt}
\setlength{\heavyrulewidth}{0.4pt}
\setlength{\lightrulewidth}{0.4pt}
\setlength{\cmidrulewidth}{0.4pt}

\title{\textbf{Apunts Química}}
\author{\textit{Gerard Cosp Lores}}
\date{\today}

% -------------------------------------------------
% TABLE OF CONTENTS
% -------------------------------------------------

\renewcommand{\contentsname}{Índex}

% Title
\renewcommand{\cfttoctitlefont}{%
    \Huge\bfseries\color{Black}
}

% Line underneath title
\renewcommand{\cftaftertoctitle}{%
    \par\vspace{0.3cm}
    \noindent\color{Black}\rule{\textwidth}{1.2pt}
    \vspace{0.4cm}
}

% Section appearance
\renewcommand{\cftsecfont}{%
    \large\bfseries\color{Black}
}

\renewcommand{\cftsecpagefont}{%
    \large\bfseries\color{Black}
}

% Subsection appearance
\renewcommand{\cftsubsecfont}{%
    \normalsize
}

\renewcommand{\cftsubsecpagefont}{%
    \normalsize\color{Black}
}

% Spacing
\setlength{\cftbeforesecskip}{10pt}
\setlength{\cftbeforesubsecskip}{3pt}

% Dotted leaders
\renewcommand{\cftsecleader}{%
    \color{Black}\cftdotfill{\cftdotsep}
}

\renewcommand{\cftsubsecleader}{%
    \color{Black}\cftdotfill{\cftdotsep}
}

% Number spacing
\setlength{\cftsecnumwidth}{2.5em}
\setlength{\cftsubsecnumwidth}{3.5em}

% -------------------------------------------------

\begin{document}
\justifying

\maketitle
\tableofcontents

\newpage

\section{Introducció als Elements}

En quimica podem dir que tenim \textbf{3 estats de la matèria}, aquests varien amb la temperatura i per tant amb l'energia cinètica que tenen.

\begin{itemize}
    \item \textbf{Gasos}, $\vec{F}_{intermoleculars} < E_c$
    \item \textbf{Líquids}, $\vec{F}_{intermoleculars} = E_c$
    \item \textbf{Sólids}, $\vec{F}_{intermoleculars} > E_c$
\end{itemize}

Definim un element quimic

\[
{}^{A}_{Z}E^{\pm x}
\]

\begin{itemize}
    \item A = Nombre de protons + Neutrons (\textit{màssic})
    \item Z = Nombre de protons (\textit{atòmic})
    \item X = Càrrega (\textit{p - e})
\end{itemize}

Quan A canvia anomenem al compost un \textbf{isòtop} i és degut a la diferència de neutrons. Per exemple l'hidrogen té 3 isòtops;

Proti; ${}^{1}_{1}H$
Deuteri; ${}^{2}_{1}H$
Triti; ${}^{3}_{1}H$

Per calcular les masses dels àtoms ho fem experimentalment. Internacionalment s'ha fixat que una \textbf{amu} és una dotzena part del que pesa l'àtom de Carboni 12. Però si ho fem experimentalment surt diferent, perquè?
Doncs degut al \textbf{defecte de massa d'Einstein} que ens diu que al perdre massa es genera una energia proporcional al quadrat de la velocitat de la llum;

$\Delta E = \Delta m c^2$

Teòricament la massa d'un àtom és la \textbf{mitjana ponderada} de tots els seus isòtops en funció de l'abundància.

\section{Orbitals Atòmics}

Si volem descriure un electró d'un àtom, fem servir l'equació d'Shrödinger de manera $\Psi(x,y,z)$ però d'aquesta manera és super tediòs trobar un electró, per això fem servir les \textbf{cordenades polars} de manera que la nostra funció queda $\Psi (r, \theta, \varphi)$, on;

\[x = r \sin \theta \cos \varphi\]
\[x = r \sin \theta \sin \varphi\]
\[z = r \cos \theta \]

Amb aixó podem obtenir unes funcións que depenen o de la distància r o dels angles de manera $R(r)Y(\theta, \varphi)$ on R depèn dels nombres quàntics n i l, i Y depèn de l i m.

L'equació d'Shrödinger té infinites solucions (\textit{comprovat matemàticament}) per arribar a una solució finita li afegim uns paràmetres que anomenem \textbf{nombres quàtnics}.

\begin{itemize}
    \item \textbf{n}, nombre primer: Aquest indica la distància entre l'electró i el nucli
    \item \textbf{l}, moment angular orbital: Aquest indica la forma del orbital (0--$n-1$) i $\{0\ (s), 1\ (p), 2\ (d), 3\ (f) \dots\}$
    \item \textbf{m}, nombre magnètic: Aquest indica l'orientació del orbital (-l - l)
    \item \textbf{s}, spin: Indica com gira l'electró ($\pm \frac{1}{2}$).
\end{itemize}

Anomenem \textbf{especie hidrogenoide} a tots aquells àtoms que tenen i \textbf{$e^-$} ($H$, $He^+$, $Li^{2+}$, $Be^{3+}$). Com que l'unica força elèctrica que existeix no te més interaccions podem definir l'energia en funció del orbital de manera;

\[ \boxed{E = \frac{-R_H}{n^2}} (R_H = 2,278\cdot 10^{-18} J = 13,6 eV)\]

I aquesta energia va en negatiu sempre ja que és una força d'atracció.

\end{document}